\documentclass[10pt,a4paper]{amsart}
\usepackage{amsmath,amssymb,amsthm,graphicx}
% Retained preamble for the Stacks Project chapter {CH} (shared project preamble).
% The chapter file pinned-original-source.tex only carries \input{{preamble}} on its first line,
% so no preamble text lives inside the chapter itself.  The Stacks Project uses one single shared
% preamble for every chapter: the snapshot stages it three times (sources/stacks-chow/preamble.tex,
% sources/stacks-etale-cohomology/preamble.tex and focused-seeds/AGP/sources/stacks-snapshot-20261011/
% preamble.tex) and those copies are identical in size.  The block below is copied mechanically,
% line for line, from sources/stacks-chow/preamble.tex lines 14-18 and 154-217: the Xy-pic setup and
% the native theorem declarations together with the mathematical macros.  Nothing was added, removed
% or redefined; the wrapper class and the amsmath/amssymb/amsthm/graphicx base packages are supplied
% by the assembler.
% For commutative diagrams we use Xy-pic
% Chapter: stacks-spaces-more-morphisms.tex -- More on Morphisms of Spaces.
\usepackage[all]{xy}

% We use 2cell for 2-commutative diagrams.
\xyoption{2cell}

% Theorem environments.
%
\theoremstyle{plain}
\newtheorem{theorem}[subsection]{Theorem}
\newtheorem{proposition}[subsection]{Proposition}
\newtheorem{lemma}[subsection]{Lemma}

\theoremstyle{definition}
\newtheorem{definition}[subsection]{Definition}
\newtheorem{example}[subsection]{Example}
\newtheorem{exercise}[subsection]{Exercise}
\newtheorem{situation}[subsection]{Situation}

\theoremstyle{remark}
\newtheorem{remark}[subsection]{Remark}
\newtheorem{remarks}[subsection]{Remarks}

\numberwithin{equation}{subsection}

% Macros
%
\def\lim{\mathop{\mathrm{lim}}\nolimits}
\def\colim{\mathop{\mathrm{colim}}\nolimits}
\def\Spec{\mathop{\mathrm{Spec}}}
\def\Hom{\mathop{\mathrm{Hom}}\nolimits}
\def\Ext{\mathop{\mathrm{Ext}}\nolimits}
\def\SheafHom{\mathop{\mathcal{H}\!\mathit{om}}\nolimits}
\def\SheafExt{\mathop{\mathcal{E}\!\mathit{xt}}\nolimits}
\def\Sch{\mathit{Sch}}
\def\Mor{\mathop{\mathrm{Mor}}\nolimits}
\def\Ob{\mathop{\mathrm{Ob}}\nolimits}
\def\Sh{\mathop{\mathit{Sh}}\nolimits}
\def\NL{\mathop{N\!L}\nolimits}
\def\CH{\mathop{\mathrm{CH}}\nolimits}
\def\proetale{{pro\text{-}\acute{e}tale}}
\def\etale{{\acute{e}tale}}
\def\QCoh{\mathit{QCoh}}
\def\Ker{\mathop{\mathrm{Ker}}}
\def\Im{\mathop{\mathrm{Im}}}
\def\Coker{\mathop{\mathrm{Coker}}}
\def\Coim{\mathop{\mathrm{Coim}}}

% Boxtimes
%
\DeclareMathSymbol{\boxtimes}{\mathbin}{AMSa}{"02}

%
% Macros for moduli stacks/spaces
%
\def\QCohstack{\mathcal{QC}\!\mathit{oh}}
\def\Cohstack{\mathcal{C}\!\mathit{oh}}
\def\Spacesstack{\mathcal{S}\!\mathit{paces}}
\def\Quotfunctor{\mathrm{Quot}}
\def\Hilbfunctor{\mathrm{Hilb}}
\def\Curvesstack{\mathcal{C}\!\mathit{urves}}
\def\Polarizedstack{\mathcal{P}\!\mathit{olarized}}
\def\Complexesstack{\mathcal{C}\!\mathit{omplexes}}
% \Pic is the operator that assigns to X its picard group, usage \Pic(X)
% \Picardstack_{X/B} denotes the Picard stack of X over B
% \Picardfunctor_{X/B} denotes the Picard functor of X over B
\def\Pic{\mathop{\mathrm{Pic}}\nolimits}
\def\Picardstack{\mathcal{P}\!\mathit{ic}}
\def\Picardfunctor{\mathrm{Pic}}
\def\Deformationcategory{\mathcal{D}\!\mathit{ef}}
\newenvironment{reference}{}{}
\newenvironment{slogan}{}{}
\newenvironment{history}{}{}
\usepackage{cleveref}
\crefname{originallemma}{original lemma}{original lemmas}
\Crefname{originallemma}{Original lemma}{Original lemmas}
\makeatletter
\newlabel{more-morphisms-lemma-flat-finite-presentation-pseudo-coherent}{{0695}{source}{}{sourcepointer-1}{}}
\newlabel{more-morphisms-lemma-flat-finite-presentation-pseudo-coherent@cref}{{[originallemma][0][]0695}{[1][0][]source}{}{}{}}
\makeatother
\crefname{originallemma}{original lemma}{original lemmas}
\Crefname{originallemma}{Original lemma}{Original lemmas}
\makeatletter
\newlabel{more-morphisms-lemma-pull-relative-pseudo-coherent}{{09US}{source}{}{sourcepointer-2}{}}
\newlabel{more-morphisms-lemma-pull-relative-pseudo-coherent@cref}{{[originallemma][0][]09US}{[1][0][]source}{}{}{}}
\makeatother
\crefname{originallemma}{original lemma}{original lemmas}
\Crefname{originallemma}{Original lemma}{Original lemmas}
\makeatletter
\newlabel{more-morphisms-lemma-qcoh-relative-pseudo-coherence-characterize}{{09VF}{source}{}{sourcepointer-3}{}}
\newlabel{more-morphisms-lemma-qcoh-relative-pseudo-coherence-characterize@cref}{{[originallemma][0][]09VF}{[1][0][]source}{}{}{}}
\makeatother
\crefname{originallemma}{original lemma}{original lemmas}
\Crefname{originallemma}{Original lemma}{Original lemmas}
\makeatletter
\newlabel{more-morphisms-lemma-relative-pseudo-coherent-descends-fppf}{{0CSN}{source}{}{sourcepointer-4}{}}
\newlabel{more-morphisms-lemma-relative-pseudo-coherent-descends-fppf@cref}{{[originallemma][0][]0CSN}{[1][0][]source}{}{}{}}
\makeatother
\crefname{originallemma}{original lemma}{original lemmas}
\Crefname{originallemma}{Original lemma}{Original lemmas}
\makeatletter
\newlabel{more-morphisms-lemma-relative-pseudo-coherent-post-compose}{{0CSP}{source}{}{sourcepointer-5}{}}
\newlabel{more-morphisms-lemma-relative-pseudo-coherent-post-compose@cref}{{[originallemma][0][]0CSP}{[1][0][]source}{}{}{}}
\makeatother
\crefname{originallemma}{original lemma}{original lemmas}
\Crefname{originallemma}{Original lemma}{Original lemmas}
\makeatletter
\newlabel{remark-match-relative-pseudo-coherence}{{0CSW}{source}{}{sourcepointer-6}{}}
\newlabel{remark-match-relative-pseudo-coherence@cref}{{[originallemma][0][]0CSW}{[1][0][]source}{}{}{}}
\makeatother
\title{Three equivalent chart tests for relative pseudo-coherence of complexes on algebraic spaces}
\author{The Stacks Project Authors}
\date{stacks-project:spaces-more-morphisms.tex:local-snapshot-20261011; GNU FDL 1.2 or later}
\begin{document}
\maketitle
\noindent\textit{Editorial note.} Statement span is lines 10761-10784 of sources/stacks-spaces-more-morphisms/pinned-original-source.tex; proof span is lines 10786-10831, the complete author proof (44 body lines).\par\medskip
\noindent\textit{Editorial note.} Navigation stubs: 6 cited-but-not-included labels are supplied as non-load-bearing pointers with their pinned tags where the snapshot tag index has them.\par\medskip
\noindent\textit{Editorial note.} extra\_packages is empty on purpose: the retained preamble already loads xy with options and a second declaration clashes in the assembler wrapper.\par\medskip
\noindent\textit{Editorial note.} Text is retained verbatim, including the author's own typos and placeholders.\par\medskip
\noindent\textit{Editorial note.} Counted B side (both complete). Central glue card for the spaces-side relative-pc toolkit of waves 12-18.\par\medskip
\noindent\textit{Allowed standard prerequisites (editorial).}\begin{itemize}
\item Explicit prerequisite boundary. More on Morphisms dictionary: qcoh-relative-pseudo-coherence-characterize, relative-pseudo-coherent-descends-fppf (L16106), relative-pseudo-coherent-post-compose (L16168), pull-relative-pseudo-coherent, flat-finite-presentation-pseudo-coherent.
\end{itemize}\medskip
\ifdefined\hypertarget\hypertarget{sourcepointer-1}{}\fi
\noindent\textit{Source reference.} Original lemma 0695: sources/stacks-spaces-more-morphisms/pinned-original-source.tex, source lines 10831--10831. Cited by the retained statement or proof but not part of the collected spans; supplied only as a navigation pointer. The argument does not depend on its content.\par\medskip
\ifdefined\hypertarget\hypertarget{sourcepointer-2}{}\fi
\noindent\textit{Source reference.} Original lemma 09US: sources/stacks-spaces-more-morphisms/pinned-original-source.tex, source lines 10831--10831. Cited by the retained statement or proof but not part of the collected spans; supplied only as a navigation pointer. The argument does not depend on its content.\par\medskip
\ifdefined\hypertarget\hypertarget{sourcepointer-3}{}\fi
\noindent\textit{Source reference.} Original lemma 09VF: sources/stacks-spaces-more-morphisms/pinned-original-source.tex, source lines 10831--10831. Cited by the retained statement or proof but not part of the collected spans; supplied only as a navigation pointer. The argument does not depend on its content.\par\medskip
\ifdefined\hypertarget\hypertarget{sourcepointer-4}{}\fi
\noindent\textit{Source reference.} Original lemma 0CSN: sources/stacks-spaces-more-morphisms/pinned-original-source.tex, source lines 10831--10831. Cited by the retained statement or proof but not part of the collected spans; supplied only as a navigation pointer. The argument does not depend on its content.\par\medskip
\ifdefined\hypertarget\hypertarget{sourcepointer-5}{}\fi
\noindent\textit{Source reference.} Original lemma 0CSP: sources/stacks-spaces-more-morphisms/pinned-original-source.tex, source lines 10831--10831. Cited by the retained statement or proof but not part of the collected spans; supplied only as a navigation pointer. The argument does not depend on its content.\par\medskip
\ifdefined\hypertarget\hypertarget{sourcepointer-6}{}\fi
\noindent\textit{Source reference.} Original lemma 0CSW: sources/stacks-spaces-more-morphisms/pinned-original-source.tex, source lines 10746--10759. Cited by the retained statement or proof but not part of the collected spans; supplied only as a navigation pointer. The argument does not depend on its content.\par\medskip
\begin{lemma}
\label{lemma-qcoh-relative-pseudo-coherence-characterize}
Let $S$ be a scheme. Let $f : X \to Y$ be a morphism of algebraic spaces
over $S$ which is locally of finite type. Let $m \in \mathbf{Z}$.
Let $E \in D_\QCoh(\mathcal{O}_X)$. With notation as explained in
Remark \ref{remark-match-relative-pseudo-coherence}
the following are equivalent:
\begin{enumerate}
\item for every commutative diagram
$$
\xymatrix{
U \ar[d] \ar[r] & V \ar[d] \\
X \ar[r] & Y
}
$$
where $U$, $V$ are schemes and the vertical arrows are \'etale, the complex
$E|_U$ is $m$-pseudo-coherent relative to $V$,
\item for some commutative diagram as in (1) with $U \to X$
surjective, the complex $E|_U$ is $m$-pseudo-coherent relative to $V$,
\item for every commutative diagram as in (1) with $U$ and $V$
affine the complex $R\Gamma(U, E)$ of $\mathcal{O}_X(U)$-modules
is $m$-pseudo-coherent relative to $\mathcal{O}_Y(V)$.
\end{enumerate}
\end{lemma}

\begin{proof}
Part (1) implies (3) by More on Morphisms, Lemma
\ref{more-morphisms-lemma-qcoh-relative-pseudo-coherence-characterize}.

\medskip\noindent
Assume (3). Pick any commutative diagram as in (1) with $U \to X$ surjective.
Choose an affine open covering $V = \bigcup V_j$ and affine open coverings
$(U \to V)^{-1}(V_j) = \bigcup U_{ij}$. By (3) and More on Morphisms, Lemma
\ref{more-morphisms-lemma-qcoh-relative-pseudo-coherence-characterize}
we see that $E|_U$ is $m$-pseudo-coherent relative to $V$.
Thus (3) implies (2).

\medskip\noindent
Assume (2). Choose a commutative diagram
$$
\xymatrix{
U \ar[d] \ar[r] & V \ar[d] \\
X \ar[r] & Y
}
$$
where $U$, $V$ are schemes, the vertical arrows are \'etale, the
morphism $U \to X$ is surjective, and $E|_U$ is $m$-pseudo-coherent
relative to $V$. Next, suppose given a second commutative diagram
$$
\xymatrix{
U' \ar[d] \ar[r] & V' \ar[d] \\
X \ar[r] & Y
}
$$
with \'etale vertical arrows and $U', V'$ schemes. We want to show
that $E|_{U'}$ is $m$-pseudo-coherent relative to $V'$.
The morphism $U'' = U \times_X U' \to U'$ is surjective \'etale
and $U'' \to V'$ factors through $V'' = V' \times_Y V$ which
is \'etale over $V'$. Hence it suffices to show that $E|_{U''}$
is $m$-pseudo-coherent relative to $V''$, see
More on Morphisms, Lemmas
\ref{more-morphisms-lemma-relative-pseudo-coherent-descends-fppf} and
\ref{more-morphisms-lemma-relative-pseudo-coherent-post-compose}.
Using the second lemma once more it suffices to show that
$E|_{U''}$ is $m$-pseudo-coherent relative to $V$.
This is true by More on Morphisms, Lemma
\ref{more-morphisms-lemma-pull-relative-pseudo-coherent}
and the fact that an \'etale morphism of schemes is pseudo-coherent by
More on Morphisms, Lemma
\ref{more-morphisms-lemma-flat-finite-presentation-pseudo-coherent}.
\end{proof}

\end{document}
